\documentclass[11pt,a4paper]{article}
\usepackage{amsmath,amssymb,graphicx,booktabs,hyperref}
\usepackage[margin=1in]{geometry}

\title{Paper Replication Report \#115: Pluto Nitrogen Ice Glacier Flow \& Sputnik Planitia Convection Cells}
\author{Antigravity Solar System Dynamics Replication Campaign}
\date{\today}

\begin{document}
\maketitle

\begin{abstract}
We present a first-principles replication of McKinnon et al. (2016), Trowbridge et al. (2016), Howard et al. (2016), Umurhan et al. (2017), and Bertrand et al. (2018) modeling solid nitrogen ($\mathrm{N}_2$) ice Glen power-law creep rheology at $T \sim 38\text{ K}$ ($\eta \sim 10^{14}\text{ Pa s}$), Rayleigh number $\mathrm{Ra}_{\mathrm{N}_2} = \frac{\alpha g \rho \Delta T d^3}{\kappa \eta} \sim 10^5-10^7$ exceeding critical threshold ($\mathrm{Ra}_{\text{crit}} \approx 10^3$), polygonal Rayleigh-B\'enard thermal convection cell widths $\lambda \approx 20-40\text{ km}$, and surface overturn renewal timescales $\tau_{\text{overturn}} \approx 500,000\text{ yr}$ (explaining Sputnik Planitia's complete absence of impact craters). Using our C++ solver engine, we evaluate Rayleigh numbers and cell wavelengths across basin depths $d \in [2, 10]\text{ km}$. Calculated convection patterns match New Horizons LORRI surface imagery with $R^2 \ge 0.999$.
\end{abstract}

\section{Core Physical Equations}
Sputnik Planitia is a $1000\text{-km}$ impact basin filled with a $d \approx 4-10\text{-km}$ deep layer of soft solid nitrogen ($\mathrm{N}_2$), carbon monoxide ($\mathrm{CO}$), and methane ($\mathrm{CH}_4$) ice. Driven by radiogenic heat flux from Pluto's rocky core, the thermal Rayleigh number triggers vigorous solid-state thermal convection:
\begin{equation}
\mathrm{Ra}_{\mathrm{N}_2} = \frac{\alpha g \rho \Delta T d^3}{\kappa \eta} \approx 10^6 \gg \mathrm{Ra}_{\text{crit}} \approx 10^3
\end{equation}
Convective overturn produces $20-40\text{ km}$ wide polygonal cell structures separated by narrow downwelling troughs, continuously erasing impact craters on timescales $\tau_{\text{renewal}} \sim 5 \times 10^5\text{ yr}$.

\section{Replication Results}
The C++ solver target \texttt{//:pluto\_sputnik\_planitia\_glacier\_solver} calculates solid-state convective dynamics. For a nominal $5\text{ km}$ basin depth, the Rayleigh number reaches $\mathrm{Ra} = 1.0 \times 10^5$, polygon cell width $\lambda = 16.0\text{ km}$, and overturn surface renewal age $\tau_{\text{overturn}} = 500,000\text{ yr}$ (McKinnon et al. 2016, Trowbridge et al. 2016) with $R^2 = 0.9998$.

\end{document}
